% -----------------------------------------------------------
\section{1864}
% -----------------------------------------------------------

	% ----------------------
	\subsection{Brigham Young}
	% ----------------------
		\label{EMS-1864-BY}

		% ----------------------
		\subsubsection{Introduction to 1864}
		% ----------------------
		
			sdf
			
		% -----
		\subsubsection{Discourse, 31 July 1864}
		% -----
			\label{EMS-1864-BY-31-JUL-1864}
			% \label{EMS-18XX-BY-DAY-3LETTERMONTH-YEAR}
			
			\begin{description}
					\item[Print Sources]
				\end{description}

					\begin{tabular}{p{0.5in}p{1in}p{0.5in}p{1in}}
					CDBY	& 4:2211--2218		& JD 		& 10:318--327 \\
					\end{tabular}
				
				\begin{description}
					\item[Online Access] \href{https://drive.google.com/file/d/1bfdBkrQaFzyVpdWGI_um2Gfm8AYGat3p/view?usp=sharing}{Here}
					% \item[Analysis] \hyperref[XX]{Here}
				\end{description}
				
				\begin{description}
					\item[Background]
				\end{description}
				
					\phantom{.}
				
				\begin{description}
					\item[Original Source]
				\end{description}
				
					Church History Library, CR 100 912 (Church History Department Pitman Shorthand transcriptions, 2013--2020), Brigham Young, 31 July 1864, transcribed by LaJean Carruth; view online at the CHL \href{https://catalog.churchofjesuschrist.org/assets/efaf279f-2278-4c31-a138-b7522c9055cc/0/0}{here} (this sermon appears in \textit{Journal of Discourses}, but with significant alterations, including in the passage quoted below)
					% brief description of original source material, with reference to JSP or somewhere else for more info
					% Background material: it might be helpful to have a note that says whether a given document was presented as a speech, was written in a journal, etc.
					% Also a comment here about how reliable the source might be, or what shape it is in, if there are large omissions or parts that are hard to understand, and so on.
					% Also a comment about when I transcribed it, whether I've checked it more than once, and so on
				
				\begin{description}
					\item[Text]
				\end{description}
				
					...our religion is the priesthood of the Son of God incorporated within the pale%
						\footnote{%
							\label{EMS-1864-BY-31-JUL-1864-pale}%
							See the 1828 dictionary for ``pale'':
								\begin{compactitem}
									\item[3] An inclosure; properly, that which incloses, like fence, limit; hence, the space inclosed. He was born within the \textit{pale} of the church; within the \textit{pale} of christianity.
									\item[4] District; limited territory.
								\end{compactitem}
							}
					of this holy priesthood we frequently hear people inquire what the priesthood is it is a pure holy system of government the priesthood is the law that governs and controls the heavenly bodies the priesthood is the power and the principle that holds all things in existence by the word of him who is the author of all things this priesthood is made manifest in the last days the Lord saw fit to speak from the heavens reveal his will to the children of men it is well known without my taking the time to notify this congregation that the Lord spoke to Joseph Smith it is well known and understood that the Latter-day Saints believe the revelations that the Lord gave through his servant Joseph consequently it does not require any delineation upon this to show to this congregation that we are brethren to Joseph believers in him as a prophet and that we subscribe to the doctrine God revealed through him this is the truth and where we are known this is known this priesthood which God has revealed is for the salvation of the children of man. This priesthood which is a perfect system of government is calculated to govern and control eventually the earth the inhabitants of the earth and all things pertaining to the earth well I might talk of liberty to refer to something that was said this morning that the enemy the opposer of Jesus the accuser of the brethren that we call Satan never did own the earth he never made a particle of it] we might take occasion to see and if you want to see the works of the enemy whenever you wish to know the design of the opposer of this you will see the result of his labor is to destroy the labor of the Son of God [illegible] to make preserve purify build up and bring up to his standard all his brethren beautify the earth purify it cleanse it take curse from it bring it to paradisaical state restore all things make all things glorious and when you wish to know the work of the enemy it is dismember opposed to defy one is to create preserve build up sanctify glorify the other is to to consume destroy waste away and bring to night this is the difference I take the liberty to say so much with regard to the work of the enemy and the time will come when it will be known] that he is a usurper and all governments that is on the face of this earth that is opposed to the government of the Son of God is a usurpation by and by it will be known I am well aware that many thinking men inquiring hearts and spirits and dispositions would ask if it is necessary for the government of God on the earth at the present day most assuredly never a time when more needed than now why because men do not know how to govern themselves without it would it be treason were we in any government there is upon the face of the earth to profess to believe in the Lord Jesus Christ to say that it is his right his inalienable right his indisputable right and prerogative to reign over the children of men and all things pertaining to the earth...
			
		% -----
		\subsubsection{Discourse, 7 September 1864}
		% -----
			\label{EMS-1864-BY-7-SEP-1864}
			% \label{EMS-18XX-BY-DAY-3LETTERMONTH-YEAR}
			
			\begin{description}
					\item[Print Sources]
				\end{description}

					\begin{tabular}{p{0.5in}p{1in}p{0.5in}p{1in}}
					CDBY	& 4:2220		& JD 		& --- \\
					\end{tabular}
				
				\begin{description}
					\item[Online Access] \href{https://catalog.churchofjesuschrist.org/assets/92e3ef48-c3c0-4d38-85d9-96e90d2b4fc6/0/0}{Here}
					% \item[Analysis] \hyperref[XX]{Here}
				\end{description}
				
				\begin{description}
					\item[Background]
				\end{description}
				
					\phantom{.}
				
				\begin{description}
					\item[Original Source]
				\end{description}
				
					Part of ``The Lost Sermons'' publishing project, in Church History Library, CR 100 912 (Church History Department Pitman Shorthand transcriptions, 2013--2020), Brigham Young, 7 September 1864, transcribed by LaJean Carruth
					% brief description of original source material, with reference to JSP or somewhere else for more info
					% Background material: it might be helpful to have a note that says whether a given document was presented as a speech, was written in a journal, etc.
					% Also a comment here about how reliable the source might be, or what shape it is in, if there are large omissions or parts that are hard to understand, and so on.
					% Also a comment about when I transcribed it, whether I've checked it more than once, and so on
				
				\begin{description}
					\item[Text]
				\end{description}
					
					I say to these Latter-day Saints some things I can tell you some things I can't tell you those things that you can't understand and will not use prudently I had better keep them myself the saints have got to be taught according to them they receive and improve upon but tell them a great many things that yet will be told the saints would be like casting pearls before swine this I don't like to do but we want the people who have received the gospel and wish to become saints indeed to receive a little here a little there and improve upon themselves little until they are perfection that is what we want and this is our labor as to the circular and what is in it it is right and still there is things in it that people don't understand [space] there is some things were not written they should understand [space] and I say to you as I can to the whole world there has not yet been a revelation given to this people from the time Joseph commenced to receive revelation but what would be [altered/ltrd?] provided the people were capable of receiving more [space] the Lord has to speak to the people according to their capacity not according to his capacity we are not prepared to receive all the heavens has for us the Lord gives revelation upon revelation here a little and there a little [space] those precepts he gives we should improve upon them as for selling of grain we take measures inasmuch as we can to have the people keep the grain still a great deal will be sold and the men that doesn't feel to sustain what the convention dictated they will sooner or later manifest their folly by going to the dark more and more more and more [sic] until they [love/lv?] the world and every project that is proposed by the leaders of this people and every plan adopted for the benefit of the people those that doesn't hear it will dwindle away sooner or later and finally will apostatize [space] I tell you some things in the circular meant to be in there that I understand you don't [space] I shall not tell it
			
% -----------------------------------------------------------
\pagebreak{}
% -----------------------------------------------------------